\usepackage{needspace}

% ============================================================
% Estrutura jurídica: artigos, cláusulas, parágrafos,
% incisos e alíneas
% ============================================================

\newcounter{artigo}
\newcounter{clausula}
\newcounter{paragrafo}
\newcounter{inciso}
\newcounter{alinea}

% ============================================================
% Medidas
% ============================================================

% Parágrafos numerados
\newlength{\recuoparagrafo}
\newlength{\larguraparagrafo}

\setlength{\recuoparagrafo}{\parindent}
\setlength{\larguraparagrafo}{2.5em}

% Parágrafo único
\newlength{\larguraparagrafounico}

\settowidth{\larguraparagrafounico}{%
  \textbf{Parágrafo único.}\enspace
}

% Incisos
\newlength{\recuoinciso}
\newlength{\largurainciso}

\setlength{\recuoinciso}{%
  \dimexpr\recuoparagrafo+\larguraparagrafo\relax
}
\setlength{\largurainciso}{3em}

% Alíneas
\newlength{\recuoalinea}
\newlength{\larguraalinea}

\setlength{\recuoalinea}{%
  \dimexpr\recuoinciso+\largurainciso\relax
}
\setlength{\larguraalinea}{2em}


% ------------------------------------------------------------
% Formatação dos números
% ------------------------------------------------------------

\renewcommand{\theartigo}{\arabic{artigo}}
\renewcommand{\theclausula}{\arabic{clausula}}
\renewcommand{\theparagrafo}{\arabic{paragrafo}}
\renewcommand{\theinciso}{\Roman{inciso}}
\renewcommand{\thealinea}{\alph{alinea}}

% ------------------------------------------------------------
% Artigo
%
% Uso:
%   \artigo
%   \artigo[Do objeto]
% ------------------------------------------------------------

\newcommand{\artigo}[1][]{%
  \par
  \Needspace{3\baselineskip}%
  \vspace{1.5em}%
  \refstepcounter{artigo}%
  \setcounter{paragrafo}{0}%
  \setcounter{inciso}{0}%
  \setcounter{alinea}{0}%
  \noindent
  \textbf{%
     Artigo~\theartigo
     \ifnum\value{artigo}<10 \textsuperscript{o}\fi%
     \if\relax\detokenize{#1}\relax
     \else\space-- #1%
     \fi
  }%
  \par\nobreak
}

% ------------------------------------------------------------
% Cláusula
%
% Uso:
%   \clausula
%   \clausula[Do objeto]
% ------------------------------------------------------------

\newcommand{\clausula}[1][]{%
  \par
  \Needspace{3\baselineskip}%
  \vspace{1.5em}%
  \refstepcounter{clausula}%
  \setcounter{paragrafo}{0}%
  \setcounter{inciso}{0}%
  \setcounter{alinea}{0}%
  \Needspace{3\baselineskip}
  \par\vspace{1.5em}
  \noindent
  \textbf{%
    Cláusula~\theclausula
    \ifnum\value{clausula}<10 \textsuperscript{a}\fi%
    \if\relax\detokenize{#1}\relax
    \else\space-- #1%
    \fi
  }%
  \par\nobreak
}


% ------------------------------------------------------------
% Parágrafo numerado
%
% Uso:
%   \paragrafo Texto...
%
% Produz:
%   § 1º Texto...
% ------------------------------------------------------------
\setlength{\recuoparagrafo}{\parindent}
\setlength{\larguraparagrafo}{2.5em}

\newcommand{\paragrafo}[1]{%
  \refstepcounter{paragrafo}%
  \setcounter{inciso}{0}%
  \setcounter{alinea}{0}%
  \par
  \begingroup
    \setlength{\parindent}{0pt}%
    \advance\leftskip by
      \dimexpr\recuoparagrafo+\larguraparagrafo\relax
    \noindent
    \llap{%
      \makebox[\larguraparagrafo][l]{%
        \textbf{%
          §~\theparagrafo
          \ifnum\value{paragrafo}<10 º\fi
        }%
      }%
    }%
    \ignorespaces#1%
    \par
  \endgroup
}

% ------------------------------------------------------------
% Parágrafo único
%
% Uso:
%   \paragrafounico Texto...
% ------------------------------------------------------------

\newcommand{\paragrafounico}[1]{%
  \setcounter{inciso}{0}%
  \setcounter{alinea}{0}%
  \par
  \begingroup
    \setlength{\parindent}{0pt}%
    \advance\leftskip by
      \dimexpr\recuoparagrafo+\larguraparagrafounico\relax
    \noindent
    \llap{%
      \makebox[\larguraparagrafounico][l]{%
        \textbf{§ único.}%
      }%
    }%
    \ignorespaces#1%
    \par
  \endgroup
}

% ------------------------------------------------------------
% Inciso
%
% Uso:
%   \inciso Texto...
%
% Produz:
%   I – Texto...
% ------------------------------------------------------------

\newcommand{\inciso}[1]{%
  \refstepcounter{inciso}%
  \setcounter{alinea}{0}%
  \par
  \begingroup
    \setlength{\parindent}{0pt}%
    \advance\leftskip by
      \dimexpr\recuoinciso+\largurainciso\relax
    \noindent
    \llap{%
      \makebox[\largurainciso][l]{%
        \textbf{\theinciso\space--}%
      }%
    }%
    \ignorespaces#1%
    \par
  \endgroup
}

% ------------------------------------------------------------
% Alínea
%
% Uso:
%   \alinea Texto...
%
% Produz:
%   a) Texto...
% ------------------------------------------------------------

\newcommand{\alinea}[1]{%
  \refstepcounter{alinea}%
  \par
  \begingroup
    \setlength{\parindent}{0pt}%
    \advance\leftskip by
      \dimexpr\recuoalinea+\larguraalinea\relax
    \noindent
    \llap{%
      \makebox[\larguraalinea][l]{%
        \textbf{\thealinea)}%
      }%
    }%
    \ignorespaces#1%
    \par
  \endgroup
}
